% main.tex — R3 TIFS-style rebuild (initial submission, <=13 pages)
% Scientific authority: frozen artifacts + audited Chinese mother manuscript (R3).
\documentclass[journal]{IEEEtran}

\usepackage{amsmath,amssymb,amsthm}
\usepackage{graphicx}
\usepackage{booktabs}
\usepackage{array}
\usepackage{multirow}
\usepackage{makecell}
\usepackage{tabularx}
\usepackage{cite}
\usepackage[hidelinks]{hyperref}
\usepackage{url}
\usepackage{xurl}

\newtheorem{proposition}{Proposition}

\newcolumntype{L}[1]{>{\raggedright\arraybackslash}p{#1}}

\begin{document}

\title{Non-One-to-One Cross-Chain Forensic Fund-Flow Correspondence: Transport Representation and Conditional Decoding}

\author{Author~Names~\IEEEmembership{Member,~IEEE}%
\thanks{Manuscript submitted for review. Author names, affiliations, and contact details are completed by the authors before submission (see AUTHOR\_METADATA\_REQUIRED.md in the submission package).}%
}

\markboth{IEEE Transactions on Information Forensics and Security}%
{Author \MakeLowercase{\textit{et al.}}: Non-One-to-One Cross-Chain Forensic Fund-Flow Correspondence}

\maketitle

\begin{abstract}
Cross-chain money laundering splits suspicious asset trails across chains, tokens, and delayed settlements, so anti-money-laundering (AML) forensics can no longer assume a one-to-one mapping between source-chain and destination-chain transactions. We formalize Cross-Chain Suspicious Fund-Flow Correspondence (CSFFC), in which fund flows are matched through a nonnegative transport plan that expresses one-to-one, one-to-many (fan-out), many-to-one (merge), and unmatched mass. We further show that ranking correspondences directly from the raw transport plan is distorted: the plan's solver dual scalings are mixed with the true correspondence affinities. To repair this, we propose directional conditional decoding---row scores normalized by column mass and column scores normalized by row mass, decoded with a mutual top-$k$ rule---whose dual-cancellation property removes the conditioned side's common scaling factor exactly; the identity is stated as a direct corollary of standard Sinkhorn structure, not as a new normalization. In a preregistered, hash-locked, one-shot confirmatory holdout on untouched seeds across three bridges (Celer, Multichain, PolyNetwork), conditional decoding raises macro edge F1 from \textbf{0.2350} (raw-plan decoding) to \textbf{0.3104}, with primary effect \textbf{$\Delta = +0.075413$} and paired 95\% CI \textbf{$[0.070557, 0.080312]$}, positive on every bridge, all five preregistered mechanism directions satisfied, and an independent verifier recomputing every result; the same experiment shows balanced and unbalanced conditional decoding differ by $-0.0003$ (CI includes zero), so no unbalanced-relaxation-specific gain is claimed. Structural experiments show that the transport representation expresses $1{\to}N$ edge sets that one-to-one decoders cannot (edge-inclusion recall 0.946/0.967, not exact topology recovery), while we report transparently that strict exact-topology recovery is zero for every tested method and a globally calibrated threshold rule is stronger on edge F1. An independent post-development data audit (data-only; no method executed) confirms real flow-level fan-out outside the development corpus (2{,}451 Tier-A fan-out units over 359.98 days) but did not meet the preregistered adequacy criterion ($G = 7 < 8$), so no external method-performance claim is made.
\end{abstract}

\begin{IEEEkeywords}
Cross-chain forensics, anti-money laundering, fund-flow correspondence, optimal transport, conditional decoding, selective abstention.
\end{IEEEkeywords}

\IEEEpeerreviewmaketitle

\input{sections/01_introduction}
\input{sections/02_related_work}
\input{sections/03_overview}
\input{sections/04_method}
\input{sections/05_experiments}
\input{sections/06_discussion}
\input{sections/07_conclusion}

\bibliographystyle{IEEEtran}
\bibliography{references}

\end{document}
